\begin{thebibliography}{99}

\small

\bibitem{ref1} Wellmann S, et al. (2010). High copeptin concentrations in umbilical cord blood after vaginal delivery and birth acidosis. Journal of Clinical Endocrinology \& Metabolism. 95(11):5091–5096. \href{https://doi.org/10.1210/jc.2010-1331}{doi:jc.2010-1331}

\bibitem{ref2} Fill Malfertheiner S, et al. (2021). Vasopressin but Not Oxytocin Responds to Birth Stress in Infants. Frontiers in Neuroscience. 15:718056. \href{https://doi.org/10.3389/fnins.2021.718056}{doi:fnins.2021.718056}

\bibitem{ref3} Sinha T, et al. (2026). Maternal influences on infant gut microbiome and health. Nature. Published online 12 August 2026. \href{https://doi.org/10.1038/s41586-026-10922-9}{doi:s41586-026-10922-9}

\bibitem{ref4} Sgritta M, et al. (2019). Mechanisms Underlying Microbial-Mediated Changes in Social Behavior in Mouse Models of Autism Spectrum Disorder. Neuron. 101(2):246–259.e6. \href{https://pubmed.ncbi.nlm.nih.gov/30522820/}{PMID:30522820}

\bibitem{ref5} Schmitt LM, et al. (2023). Results of a phase Ib study of SB-121, an investigational probiotic formulation, a randomized controlled trial in participants with autism spectrum disorder. Scientific Reports. 13:5192. \href{https://doi.org/10.1038/s41598-023-30909-0}{doi:s41598-023-30909-0}

\bibitem{ref6} Mazzone L, et al. (2024). Precision microbial intervention improves social behavior but not autism severity: A pilot double-blind randomized placebo-controlled trial. Cell Host \& Microbe. 32(1):106–116.e6. \href{https://pubmed.ncbi.nlm.nih.gov/38113884/}{PMID:38113884}

\bibitem{ref7} Howard MC, Davis MM, Serviss E. (2026). A meta-analysis of perceived infectability, germ aversion, disgust and outgroup perceptions: Evaluating research on the behavioural immune system. British Journal of Psychology. Published online 18 April 2026. \href{https://doi.org/10.1111/bjop.70076}{doi:bjop.70076}

\bibitem{ref8} Iwasaki S, Takahashi T. (1998). Developmental changes in calcium channel types mediating synaptic transmission in rat auditory brainstem. Journal of Physiology. 509(Pt 2):419–423. \href{https://pubmed.ncbi.nlm.nih.gov/9575291/}{PMID:9575291}

\bibitem{ref9} Gorman KM, et al. (2019). Bi-allelic Loss-of-Function CACNA1B Mutations in Progressive Epilepsy-Dyskinesia. American Journal of Human Genetics. 104:948–956. \href{https://pubmed.ncbi.nlm.nih.gov/30982612/}{PMID:30982612}

\bibitem{ref10} Eroglu C, et al. (2009). Gabapentin receptor \ensuremath{\alpha_2\delta}-1 is a neuronal thrombospondin receptor responsible for excitatory CNS synaptogenesis. Cell. 139(2):380–392. \href{https://doi.org/10.1016/j.cell.2009.09.025}{doi:j.cell.2009.09.025}

\bibitem{ref11} Risher WC, et al. (2018). Thrombospondin receptor \ensuremath{\alpha_2\delta}-1 promotes synaptogenesis and spinogenesis via postsynaptic Rac1. Journal of Cell Biology. 217(10):3747–3765. \href{https://pubmed.ncbi.nlm.nih.gov/30054448/}{PMID:30054448}

\bibitem{ref12} Ferron L, Kadurin I, Dolphin AC. (2018). Proteolytic maturation of \ensuremath{\alpha_2\delta} controls the probability of synaptic vesicular release. eLife. 7:e37507. \href{https://doi.org/10.7554/eLife.37507}{doi:eLife.37507}

\bibitem{ref13} Merke DP, et al. (2013). Tenascin-X haploinsufficiency associated with Ehlers-Danlos syndrome in patients with congenital adrenal hyperplasia. Journal of Clinical Endocrinology \& Metabolism. 98(2):E379–E387. \href{https://pubmed.ncbi.nlm.nih.gov/23284009/}{PMID:23284009}

\bibitem{ref14} Weaver ICG, et al. (2004). Epigenetic programming by maternal behavior. Nature Neuroscience. 7:847–854. \href{https://doi.org/10.1038/nn1276}{doi:nn1276}

\bibitem{ref15} McGowan PO, et al. (2009). Epigenetic regulation of the glucocorticoid receptor in human brain associates with childhood abuse. Nature Neuroscience. 12:342–348. \href{https://doi.org/10.1038/nn.2270}{doi:nn.2270}

\bibitem{ref16} Barack DL, et al. (2024). Attention deficits linked with proclivity to explore while foraging. Proceedings of the Royal Society B. 291:20222584. \href{https://doi.org/10.1098/rspb.2022.2584}{doi:rspb.2022.2584}

\bibitem{ref17} Dąbrowski K. (1964). Positive Disintegration. Boston: Little, Brown. \href{https://archive.org/details/positivedisinteg00dabr}{source file}

\bibitem{ref18} Zvejniece L, et al. (2015). R-phenibut binds to the \ensuremath{\alpha}2-\ensuremath{\delta} subunit of voltage-dependent calcium channels and exerts gabapentin-like anti-nociceptive effects. Pharmacology Biochemistry and Behavior. 137:23–29. \href{https://doi.org/10.1016/j.pbb.2015.07.014}{doi:j.pbb.2015.07.014}

\bibitem{ref19} Xiao J, et al. (2026). Maternal Cesarean Section and Offspring ASD or ADHD Risk: A Nurses’ Health Study II Analysis. American Journal of Epidemiology. Online publication. \href{https://pubmed.ncbi.nlm.nih.gov/42363613/}{PMID:42363613}

\bibitem{ref20} Zhou C, et al. (2026). Effects of vaginal microbiota transfer on social-emotional and neurodevelopment in cesarean-born infants: 12-month follow-up of a pilot randomized clinical trial. Acta Obstetricia et Gynecologica Scandinavica. 105(6):1064–1074. \href{https://doi.org/10.1111/aogs.70218}{doi:aogs.70218}

\bibitem{ref21} Doshi AR, Hauser OP. (2024). Generative AI enhances individual creativity but reduces the collective diversity of novel content. Science Advances. 10(28):eadn5290. \href{https://doi.org/10.1126/sciadv.adn5290}{doi:sciadv.adn5290}

\bibitem{projP1} Ellis ET. (2026). The RCCX–Hunter–Dąbrowski Cascade, v3.2. Unpublished working framework; supplied canonical source. \href{https://github.com/ETEllis/DRC/blob/main/papers/rccx-hunter-dabrowski-v3.2.pdf}{source file}

\bibitem{projP2} Ellis E. (2026). The Polymathic Singularity, v0.1. Unpublished working paper; supplied canonical source. \href{https://github.com/ETEllis/DRC/blob/main/papers/polymathic-singularity-v0.1.pdf}{source file}

\bibitem{projP3} Ellis ET. (2026). Perinatal Regulatory Initialization and the \ensuremath{\alpha_2\delta}–CaV2 Developmental Hinge. Unpublished extension draft; supplied canonical source. \href{https://github.com/ETEllis/DRC/blob/main/papers/perinatal-extension-source.md}{source file}

\bibitem{ref22} Kim et al. (2023). Molecular basis and genetic testing strategies for diagnosing 21-hydroxylase deficiency, including CAH-X syndrome. Annals of Pediatric Endocrinology \& Metabolism. \href{https://pubmed.ncbi.nlm.nih.gov/37401054/}{PMID:37401054}

\bibitem{ref23} Sekar et al. (2016). Schizophrenia risk from complex variation of complement component 4. Nature. \href{https://pubmed.ncbi.nlm.nih.gov/26814963/}{PMID:26814963}

\bibitem{ref24} Demontis et al. (2023). Genome-wide analyses of ADHD identify 27 risk loci, refine the genetic architecture and implicate several cognitive domains. Nature Genetics. \href{https://pubmed.ncbi.nlm.nih.gov/36702997/}{PMID:36702997}

\bibitem{ref25} Ramo-Fernández et al. (2019). The effects of childhood maltreatment on epigenetic regulation of stress-response associated genes: an intergenerational approach. Scientific Reports. \href{https://pubmed.ncbi.nlm.nih.gov/31000782/}{PMID:31000782}

\bibitem{ref26} Bierer et al. (2020). Intergenerational Effects of Maternal Holocaust Exposure on FKBP5 Methylation. American Journal of Psychiatry. \href{https://pubmed.ncbi.nlm.nih.gov/32312110/}{PMID:32312110}

\bibitem{ref27} Metin et al. (2012). A meta-analytic study of event rate effects on Go/No-Go performance in attention-deficit/hyperactivity disorder. Biological Psychiatry. \href{https://pubmed.ncbi.nlm.nih.gov/23062355/}{PMID:23062355}

\bibitem{ref28} Samson et al. (2017). Chronotype variation drives night-time sentinel-like behaviour in hunter-gatherers. Proceedings of the Royal Society B. \href{https://pubmed.ncbi.nlm.nih.gov/28701566/}{PMID:28701566}

\bibitem{ref29} Nieto-Rostro et al. (2023). Nerve injury increases native CaV2.2 trafficking in dorsal root ganglion mechanoreceptors. Pain. \href{https://pubmed.ncbi.nlm.nih.gov/36524581/}{PMID:36524581}

\bibitem{ref30} Schlick et al. (2010). Voltage-activated calcium channel expression profiles in mouse brain and cultured hippocampal neurons. Neuroscience. \href{https://pubmed.ncbi.nlm.nih.gov/20188150/}{PMID:20188150}

\bibitem{ref31} El-Awaad et al. (2019). Direct, gabapentin-insensitive interaction of a soluble form of the calcium channel subunit \ensuremath{\alpha_2\delta}-1 with thrombospondin-4. Scientific Reports. \href{https://pubmed.ncbi.nlm.nih.gov/31700036/}{PMID:31700036}

\bibitem{ref32} Hull et al. (2019). Development and Validation of the Camouflaging Autistic Traits Questionnaire (CAT-Q). Journal of Autism and Developmental Disorders. \href{https://pubmed.ncbi.nlm.nih.gov/30361940/}{PMID:30361940}

\bibitem{ref33} Celletti et al. (2020). A new insight on postural tachycardia syndrome in 102 adults with hypermobile Ehlers-Danlos Syndrome/hypermobility spectrum disorder. Monaldi Archives for Chest Disease. \href{https://pubmed.ncbi.nlm.nih.gov/32434316/}{PMID:32434316}

\bibitem{ref34} Klinger-König et al. (2019). Methylation of the FKBP5 gene in association with FKBP5 genotypes, childhood maltreatment and depression. Neuropsychopharmacology. \href{https://pubmed.ncbi.nlm.nih.gov/30700816/}{PMID:30700816}

\bibitem{ref35} Dahimene et al. (2022). Biallelic CACNA2D1 loss-of-function variants cause early-onset developmental epileptic encephalopathy. Brain. \href{https://pubmed.ncbi.nlm.nih.gov/35293990/}{PMID:35293990}

\end{thebibliography}
